% !TeX root = ../papers/nc-p3.tex
\section{行列式比の円錐曲線族の全評価と三次元無印商}
\label{sec:unequal-determinant-conic}
節\ref{sec:full-thin-parallel-cover}の $u=1,v\ne1,z=-1,\lambda=1$ の表示で、
$A=\left(\begin{smallmatrix}a&b\\b&d\end{smallmatrix}\right)$,
$B=\left(\begin{smallmatrix}e&b\\b&h\end{smallmatrix}\right)$,
$D=ad-b^2$, $\beta=(eh-b^2)/D$ とする。$D\beta(1-\beta)\ne0$ である。
$e-h=\beta(a-d)$ により一意な $k$ があり、
\[
 e=k+\beta(a-k),\qquad h=k+\beta(d-k),\qquad
 \beta(a-k)(d-k)=(k-b)(k+b).
\]
最後の式は行列式比から得られ、逆も成り立つ。以下
$I=(a-b)(e-b),\ J=(a+b)(e+b),\ K=(d-b)(h-b),\ N_0=(d+b)(h+b)$、
$T_\pm=e\pm b+\beta(d\pm b)$ とする。

\begin{theorem}\label{thm:unequal-conic-all-evaluations}
四辺の全ての非零左右評価が階数3以上となる必要十分条件は
$(1+\beta)I J K N_0\ne0$ であり、その外に標準AS窓はない。
この開集合の反対階数は10、四辺は各々長さ20の有限スキームの同型のグラフである。
偶位置0の平方座標の支持は次の六点の全逆像である。
\[
\begin{gathered}
 (1,1,0,0),\quad(0,0,1,-1),\\
 (0,-2(d+h),-N_0,K),\quad(-2(a+e),0,-J,I),\\
 (-N_0,-J,0,D T_+),\quad(K,I,D T_-,0).
\end{gathered}
\]
異なる支持点数は $20-4[d+h=0]-4[a+e=0]-4[T_+=0]-4[T_-=0]$。
合流点は長さ3の曲線状局所環を持つ。正則なら最小周期は2である。
\end{theorem}
\begin{proof}
偶点 $x=(u,v,s,t)$ で奇点に作用する六式の行列は
\[
 E(x)=\begin{pmatrix}
 u&0&ms&-nt\\0&v&fs&gt\\0&u&Nt&Ms\\v&0&-Gt&Fs\\
 s&-t&u&v\\t&s&v&u
 \end{pmatrix},
\]
ここで $(m,n,f,g)=(d-b,d+b,a-b,a+b)/D$,
$(M,N,F,G)=(h-b,h+b,e-b,e+b)/(\beta D)$。
厚い二角の列基底から得られる。反対の偶位置では最後の偶座標の符号を変えれば同じである。
実際二移動の対角倍率は偶で $(-1,-1,-1,1)$、奇で恒等となり、
偶交換から逆辺は転置である。

まず $I J K N_0\ne0$ とする。$k=\pm b$ なら $a$ または $d$ が $k$ に等しくなるので、
ここでは $(a-k)(d-k)(k^2-b^2)\ne0$。行番号を0から数え、四行の最大行列式を $[ijkl]$ とすると
\[
 (m-M)[0123]+Mf[0345]-Fm[1245]=cuvst,
\quad c=\frac{-2(a-b)(a+b)(b-d)(b+d)(a-k)(d-k)(1+\beta)}
 {(b-k)(b+k)^2D^3}.
\]
これは上の行列と円錐の式の展開である。$\beta\ne-1$ なら一座標が消える。
$u=0,vst\ne0$ では核の高位座標を $(nt,ms)$ に取れ、
第三式で $mM s^2+nN t^2=0$、第二・第四式で二低位座標が決まる。
薄い二式は $mM v^2=-2(d+h)t^2/(\beta D^2)$ に還元される。
$v=0,ust\ne0$ では同じく高位を $(gt,-fs)$ に取ると第四表示を得る。
$s=0,uvt\ne0$ では高位座標の比は $y_2/y_3=nv/(Gu)=gu/(Nv)$、
薄い式は $(G+n)v^2=-gG t^2$ となる。
$t=0,uvs\ne0$ では $y_3/y_2=mv/(Fu)=fu/(Mv)$ と $(M+f)u^2=mM s^2$ である。
ここでは $fn+gm=2/D$, $FN+GM=2/(\beta D)$,
$M+f=F+m$, $N+g=G+n$ を使った。各計算の核は一本である。
二座標が零なら、低位と高位に一つずつ非零のときは厚い式が核を零にする。
同じ組に二つ非零のときは表の最初の二表示となり、厚い式は他方の二座標を自由にする。
薄い式でその一本が残る。一座標のみの点では階数4である。
よって左の像は有限で全核は一本、右の点に直線の逆像があれば左の像の有限性に反する。
六つの双一次式はゼロ次元完全交叉で、その長さは $\binom63=20$ である。
反対の二移動と転置により全四辺に通用する。反対切断の四色の階数は $4,2,2,2$。
第一行列式の分子は定数倍を除いて
$((a-k)(d-k)(1-\beta))^2(a-k)(d-k)(1+\beta)$、分母はここで単元であり、
他の三行列の独立二行は直接取れる。周期は系\ref{cor:minus-one-second-completion-periods}による。

$a+e=d+h=0$ は $\beta=1$ を意味し、$T_+=T_-=0$ なら $b=0$ かつ
$e=-\beta d,h=-\beta a$ となり再び $\beta=1$ である。
従って高位と低位の各二点対には一つの四点軌道しか衝突しない。
稠密開部の異なる二十点から一般の付値曲線で特殊化し、平方根を分岐後に取ると、
各衝突点に元の一点と動く二点が来る。特殊側も長さ20なので局所長はちょうど3である
（付値環上の一般ファイバーの閉包は有限平坦）。低位点 $(1,1,0,0)$ では
核を $(0,0,1,-1)$ として一次化すると $dv=0$ と
$(M+f)ds=(G+n)dt=0$ が得られる。少なくとも一係数が非零である。
高位点 $(0,0,1,i)$, $i^2=-1$ では薄い式で $dt=0$ と低位核の比の微分が決まる。
二つの高位核座標の変分について残る二式の係数は
$(i(m+N),n+M)$ と $(-i(f+G),g+F)$ で、$u,v$ の変分はこれらの座標から決まる。
$d+h=0$ なら第一行は $(1-1/\beta)(im,n)$、$a+e=0$ なら第二行が
$(1-1/\beta)(-if,g)$ となる。他の符号でも同様なので接次元は1である。

$\beta=-1$ では $I J K N_0\ne0$ の場合、$Z^2=-mN/(nM)$,
$X^2=m(M+f)/M$ とする。$M+f=(D/k)/D$ は符号を除いて $1/k\ne0$ である
（$D=k(a+d-2k)$）。奇点 $(X,MZX/m,1,Z)$ での評価の階数は2である。
実際 $FN=-fn,GM=-gm$ で後ろの厚い二式は先頭の二式の倍数となり、
薄い二式にも同じ代入をすれば両方がその包に入る。

最後に $I J K N_0=0$ では、両角の交換、両対角の同時交換、$b$ の符号変更を使い
$a=b$ としてよい。これらは薄い標準語を保つ付け替えと倍率である（下の商の計算）。
$D\ne0$ なので $b(d-b)\ne0$。円錐の式により $e=b$ または $h=-b$。
前者では $f=F=0$, $m=M=1/b$。$x=(r,0,1,0)$, $r^2=1/b$ で行列が階数2になる。
後者では $f=N=0$, $m=-G=1/b$ なので、同じ座標の奇点で転置評価が階数2になる。
定理\ref{thm:single-opposite-evaluation-ranks}で全てが排除される。
\end{proof}

\begin{corollary}\label{cor:unequal-conic-unmarked-quotient}
一般の射影安定群は文字の四元群である。この族の一般の無印商は有理三次元である。
$b(a-d)\ne0$ の一般点で
\[
 x=(a+d)/(2b),\quad y=(e+h)/(2b),\quad z=(a-d)/(2b),\quad
 A_0=y/x,\quad B_0=\beta/A_0,\quad C_0=\beta z^2
\]
と置く。$A_0,B_0,C_0$ の各々に $t\mapsto(t-1)/(t+1)$ を施した値を $r,s,t$ とすると、
三つの無印座標は $(r+s)^2,(r-s)^2,(r+s)(r-s)t$ である。
特別な全ての評価適合点にも、独立四点の組と文字倍率による有限の同型比較が適用できる。
\end{corollary}
\begin{proof}
$(a,b,d,e,h)=(5,2,11,7,37)$ の素数761での偶配置・奇配置に対し、
上の表の平方根と一意な核を取る。両方とも相異なる二十点である。
三点の従属と平面上の点数で点を分け、独立四点と全係数非零の第五点の像を列挙すると、
射影自己写像数は偶で8、奇で4である。両方で一平面の最大点数は6であり、
任意の直線上の点数も6以下である。従って一様重みの配置は安定である。
具体的には、付値環上で異なる二十点を分離しておくと、射影写像の行列の原始な特殊化が
階数 $j<4$ なら核の外の全点が $j$ 次元線型空間の射影に入る。
核内と像内の点数の上界はそれぞれ $5(4-j)$ と $5j$ より真に小さいので不可能である。
従って一般の自己同型候補の置換もこの標本に単射で特殊化する。
全窓の自己同型は奇位置の配置に忠実に作用する。実際一点の相手は一意で、
各位置の配置が射影枠を含むからである。文字群が位数4を実現するので一般の群はこれに等しい。

従って一般点の無印同型は必ず文字群を運ぶ。四位置の四固有線の付け替えと
文字倍率だけを考えてよい。$y^2=\beta x^2+(1-\beta)(1-\beta z^2)$。
薄い二色のラベルを保ち規格化した語を移す置換を、節\ref{sec:full-thin-parallel-cover}と
同じ整数文字行列で解くと、$b$ による共通倍率除去の後の作用は16個である。生成元は
$(x,y)\mapsto(-x,-y)$, $z\mapsto-z$ に加えて
\[
 (x,y,z,\beta)\longmapsto(y,x,\beta z,1/\beta),\qquad
 (x,y,z,\beta)\longmapsto(-x/z,-y/(\beta z),1/z,1/\beta).
\]
前二者では書かなかった座標を固定する。
実際、低位と高位を動かさない置換では両角の交換、対角の同時交換と符号を得る。
全位置で高低を交換すればSchur補完が新しい上角となり、最後の式を得る。
再現計算は全位置移動と共通の文字せん断を直接列挙する。
最初の二符号の不変体は $A_0,B_0,C_0$ が生成する。実際 $y/x,\beta,z^2$ から
二次式で $x^2$ が復元できる。残る群は $(A_0,B_0,C_0)$ の初めの二数を同時に逆数にする操作と、
その二数の交換および最後の逆数を同時に取る操作である。Cayley変換の後は
$(r+s,r-s,t)$ の符号積を保つ四通りの変更なので、表示の不変体を得る。
例外値の比較には、全位置で元の配置の二つの固定点対が独立な四点を与える。
核の移送でも二つの固定対は隣の二座標直線上の二対に移る。
従って各位置で支持の独立四点を選び、前節の対角倍率の整数文字判定で四窓全体を比較できる。
これは各閉点の全次数正則性を仮定せずに同型を判定する。
\end{proof}
